% !TeX root = ../thuthesis-example.tex

\chapter{绪论（课题背景、现状）}

\section{课题背景}

电子商务平台每天会产生大量用户行为日志，包括浏览、加购、移除购物车、购买、商品、类目、品牌、价格、用户和会话等字段。单机脚本可以完成小样本统计，但当数据量从几十 KB 的演示文件增长到数 GB 的真实日志时，原有处理方式会暴露出三个问题：第一，CSV 扫描成本高，重复读取会拖慢整个分析流程；第二，复杂分析模块容易把明细结果收集到 driver 端，造成 Python list 或 JSON payload 过大；第三，关联分析、推荐候选、异常检测等模块存在组合膨胀风险，数据规模增加后可能出现长尾 task、shuffle 放大和内存溢出。

本课题的目标不是简单把程序放进 Docker，也不是仅把 Spark master 改成 YARN，而是围绕“真实大数据实验是否可复现、可解释、可对比”设计一套本机稳妥型实验系统。系统使用 Kaggle 电商行为数据作为原始输入，保留 Oct 和 Nov 两个月 CSV 原始数据，通过用户级采样构造 1\% 与 5\% 实验集，再在 Docker ARM64 Hadoop YARN 集群上运行 PySpark pipeline。最终产出包括 HDFS 数据快照、Spark 运行记录、质量报告、前端 Ops 页和 benchmark 对比结果。

\section{国内外现状}

大数据批处理的经典起点是 MapReduce。MapReduce 将海量数据处理抽象为 Map 和 Reduce 两个阶段，通过分布式文件系统、任务重试和数据本地性降低大规模批处理的工程复杂度\cite{dean2004mapreduce}。Hadoop 生态在此基础上形成了 HDFS、MapReduce 和 YARN 等组件。YARN 将资源管理从具体计算框架中抽离出来，由 ResourceManager、NodeManager 和 ApplicationMaster 协同完成资源调度，使同一集群能够承载不同计算引擎\cite{vavilapalli2013yarn}。

Spark 在 RDD 模型中引入内存计算、lineage 容错和数据复用，解决了 MapReduce 在迭代计算、交互式分析和多阶段作业中的效率问题\cite{zaharia2012rdd,zaharia2016spark}。Spark SQL 又通过 Catalyst 优化器、DataFrame API、列式执行和谓词下推，使结构化数据分析更加接近数据库优化器的工作方式\cite{armbrust2015spark}。在现代 Spark 应用中，是否启用 AQE、是否使用 Parquet、是否合理控制 shuffle 分区和 driver collect，往往比“是否使用集群”本身更影响稳定性。

同时，学术界和工业界已经认识到大数据系统性能不能只依赖静态经验参数。Starfish 提出 profile-driven tuning，强调先收集运行画像，再根据证据调优\cite{herodotou2011starfish}；SkewTune 关注数据倾斜和长尾 task 对整体作业时间的影响\cite{kwon2012skewtune}；Themis 则提醒系统设计不能盲目追求全内存，应当接受合理 spill 和磁盘中间结果\cite{rasmussen2012themis}。这些思想与本课题的实验结论一致：仅增加 YARN 不能自动避免 OOM，必须把资源调度、执行计划、数据布局和算法输出边界一起纳入设计。

\section{课题问题定义}

本项目早期使用 local Spark 和小样本 JSON 缓存，适合演示页面功能，但不足以支撑“大数据与云计算”课程中的集群实验和论文结论。前端出现红色长条警告、滚动异常等问题时，根因并不只在页面布局，也与后端一次性输出过多明细、质量状态不清晰和运行证据分散有关。因此，本课题将问题拆成四类：

\begin{enumerate}
    \item \textbf{数据规模问题}：从 87.4KB smoke 文件扩展到真实 Oct/Nov 原始 CSV，并构造 1\% 和 5\% 用户级样本。
    \item \textbf{运行环境问题}：从 \texttt{local[*]} 升级为 Docker Hadoop YARN 实验集群，同时兼容 ARM64 Mac，避免下载 amd64 Hadoop 镜像。
    \item \textbf{算法内存问题}：限制 driver 端 full collect，将明细输出收敛为 Parquet 与 preview JSON，并对关联、推荐、异常等模块加入上界。
    \item \textbf{证据展示问题}：把 benchmark、HDFS 输入、YARN app id、质量状态和清理策略汇总到前端 Ops 页与 LaTeX 报告中，减少答辩时翻终端的成本。
\end{enumerate}

\section{本文结构}

本文第 2 章介绍 Hadoop、YARN、Spark、Parquet、Docker 和前端运维展示等相关技术；第 3 章给出系统需求、设计目标和验收指标；第 4 章详细说明 Docker YARN 集群、Spark 提交流程、核心代码实现、数据清理、前端展示和 benchmark 实验结果；第 5 章为致谢；最后列出参考文献。
